\chapter{Introduction}

\section{Background and Motivation}

Epilepsy is a common neurological disorder: an estimated 45.9 million people had active epilepsy in 2016
\cite{gbd2019epilepsy}. Seizures are classified by onset; focal seizures start in one part of the brain, may or may
not impair awareness, and can spread into bilateral tonic-clonic seizures \cite{fisher2017ilae}. Patients are often
unaware of their own seizures, so seizure diaries undercount them \cite{elger2018underreporting}, and tonic-clonic
seizures carry a risk of sudden unexpected death in epilepsy \cite{ryvlin2013mortemus}. Reliable, automatic seizure
detection outside the hospital would give clinicians objective seizure counts and caregivers timely alarms.

In hospital, seizures are diagnosed with video-EEG from a full scalp electrode cap, which cannot be worn at home for
days. Wearable devices are recommended today only for detecting tonic-clonic seizures \cite{beniczky2021guideline};
for focal seizures that do not become tonic-clonic (822 of the 883 seizures in SeizeIT2) the problem is open, and
false alarms are the main obstacle \cite{swinnen2026seizeit2}. Behind-the-ear EEG combined with ECG, EMG and motion
sensors is a discreet alternative \cite{vandecasteele2020bte,vandecasteele2021ecg}, and the SeizeIT2 dataset
\cite{bhagubai2025seizeit2} makes it possible to develop and evaluate detectors on 125 patients recorded with such a
wearable.

A wearable detector has to run for days on a small battery, close to the sensors, and raise an alarm within seconds.
Sending raw signals to a phone or a server costs radio energy and depends on a connection; running a neural network on
a microcontroller or an accelerator costs multiplications, which dominate the area and power of digital hardware.
Field-programmable gate arrays (FPGAs) offer a different trade-off: their basic logic element is a small lookup table
(LUT), typically with six inputs, that implements any Boolean function of those inputs in one clock cycle.

A family of recent machine-learning models, which we call \emph{LUT-native}, is built from exactly this element:
weightless neural networks, whose neurons are truth tables \cite{bacellar2024dwn}, logic-gate networks
\cite{petersen2022dlgn}, and LUT-based networks designed for FPGAs \cite{umuroglu2020logicnets}. They need no
multipliers at inference time, so a classifier can cost a few hundred to a few thousand LUTs. They have been applied
to image benchmarks, physics triggers and, very recently, ECG \cite{mommen2026ecg} and EEG \cite{dharia2026eeglgn}
classification, but, to our knowledge, not to seizure detection on multimodal wearable data.

\section{Problem Statement and Objectives}

\begin{quote}
\itshape Can a LUT-native neural network, fed by integer features computed from behind-the-ear EEG, ECG, EMG and
motion, detect seizures in patients it has never seen as well as the published SeizeIT2 baselines, with no
multipliers in the classifier and a budget of a few thousand FPGA LUTs?
\end{quote}

The objectives of the project are:
\begin{enumerate}
\item Audit the SeizeIT2 dataset \cite{bhagubai2025seizeit2} completely and reproduce its published baselines.
\item Design an integer-only feature pipeline that a hardware front end can compute bit-exactly, and show that it loses
  no accuracy against floating point.
\item Compare LUT-native models (DWN, logic-gate networks, TreeLUT, LogicNets-style networks) with software
  references (gradient-boosted trees) and multiplier-based networks (INT8 MLP, CNN/TCN, GRU) on the same features,
  splits and metrics, reporting accuracy against estimated hardware cost.
\item Extend the best LUT-native model with ideas specific to wearable seizure detection (memory across seconds,
  cheaper input encoders, adaptation to new patients, learning from a stronger model, losses for rare and noisily
  labelled events, robustness to missing sensors, interpretability), keeping only what is statistically significant.
\item Evaluate the final models once on the 29 held-out test patients with event-based metrics (SzCORE
  \cite{dan2024szcore}).
\item Deliver a hardware handoff package (truth tables, thresholds, test vectors, verified reference Verilog) and
  implement the detector on an FPGA.
\end{enumerate}

Objectives 1--3 and the first part of objective 4 form Phase I and are reported here; objectives 4--6 are completed in
Phase II.

\section{Expected Contributions}

\begin{itemize}
\item The first differentiable LUT-native seizure detector, and the first benchmark of LUT-family models on SeizeIT2.
\item A bit-exact integer multimodal feature front end with its fixed-point specification.
\item An independent reproduction of the SeizeIT2 ChronoNet baseline, including a protocol that uses no test
  information.
\item Extensions of the DWN for biosignal events (stateful LUT networks, comparator-budgeted encoders, robustness to
  missing sensors), each judged by per-patient statistics.
\item A hardware handoff package that lets the FPGA implementation be checked bit for bit against the software.
\end{itemize}

\section{Organisation of the Report}

Chapter~\ref{ch:lit} reviews the literature on wearable seizure detection, detection algorithms and their evaluation,
LUT-native networks and hardware seizure detectors, and identifies the research gaps. Chapter~\ref{ch:method}
describes the dataset and our audit of it, the proposed system, the models compared and the evaluation protocol.
Chapter~\ref{ch:results} presents the experimental results of Phase I. Chapter~\ref{ch:challenges} describes the
problems we met and how we solved them. Chapter~\ref{ch:concl} discusses the findings, summarises the progress and
gives the plan for Phase II. Appendix~\ref{app:sources} lists the source file of every reported number.
